\documentclass[10pt,a4paper]{amsart}
\usepackage[margin=19mm]{geometry}
\usepackage{amsmath,amssymb,mathtools}
\usepackage[scr=rsfs,cal=euler]{mathalfa}
\usepackage{xfrac}
\usepackage[unicode,hidelinks,bookmarks=false]{hyperref}
\newcommand{\ie}{i.e.\ }
\newcommand{\cf}{cf.\ }
\newcommand{\resp}{respectively\ }
\newcommand{\Subsection}{\subsection}
\providecommand{\nobreakdash}{\nobreak\hbox{-}}
\setlength{\emergencystretch}{2em}
\multlinegap0pt
\usepackage[shortlabels]{enumitem}
\usepackage{amssymb}
\usepackage{mathtools}
\usepackage{booktabs}
\usepackage{tikz}
\usepackage{bm}
\usepackage{bbm}
\usepackage{dsfont}
\usepackage{braket}
\usepackage{caption}
\usepackage{cancel}
\usepackage{subfig}
\newtheorem{theorem}{Theorem}
\newtheorem{prp}{Proposition}
\newtheorem{cor}{Corollary}
\newcommand{\R}{\numberset{R}}
\newcommand{\numberset}{\mathbb}
\title{Sharp non-explicit blow-up profile for perturbed nonlinear heat equations with gradient terms}
\author{Maiss\^a Boughrara}
\date{Source: arXiv:2310.06620v2 (CC BY 4.0)}
\begin{document}
\maketitle

\noindent\emph{Source and licence.} Sharp non-explicit blow-up profile for perturbed nonlinear heat equations with gradient terms. Version arXiv:2310.06620v2 (CC BY 4.0), distributed by arXiv under the Creative Commons Attribution 4.0 International licence, http://creativecommons.org/licenses/by/4.0/, as recorded in the arXiv licence metadata for this exact version and shown on the abstract page https://arxiv.org/abs/2310.06620v2. Full licence notice: \texttt{licenses/arxiv-2310.06620-CC-BY-4.0.txt}.\par\smallskip

\noindent\textit{Editorial scope.} Counted unit: the article's theorem g (source0.tex lines 141--150). The article's complete original proof of it is delivered with it (source0.tex lines 1991--2157, one \texttt{proof} environment of 167 lines, which the article places away from the statement). No mathematics is written, repaired or re-derived: the statement and the proof are the article's own lines, in the article's own notation, and no retained line is substituted or re-flowed.  Retained uncounted as the source-local context the proof needs: lines 92--97; lines 99--101; lines 113--115; lines 126--128; lines 1270--1272; lines 1274--1276; lines 1311--1316; lines 1341--1346; lines 1356--1358; lines 1391--1419; lines 1416--1418; lines 1844--1846; lines 1923--1925.  Nothing was omitted from any retained span, so the package carries no omission record.  No retained line contains a citation, so the wrapper carries no bibliography and no bibliography entry.  No retained line contains a figure environment, an image-inclusion command or drawing code, so no image is delivered and the package declares no asset dependency.  Editorial formatting only, in addition: an AMS \texttt{amsart} preamble that restates the article's own declarations and macros used by the retained lines, namely the environments \texttt{theorem}, \texttt{prp}, \texttt{cor} on separate counters, together with the standard packages those lines need.  The retained environments print in this document as Theorem 1, Prp 1, Corollary 1 in order of appearance, whereas the article prints the same items with the numbering of its own full document.  The complete native source of the article as distributed in the arXiv e-print for this exact version is bundled unchanged as \texttt{sources/148262/source0.tex}.
\begin{equation}\label{eq:1}
\begin{split}
u_t=\Delta u+|u|^{p-1}u+h(u,\nabla u),\\
u(.,0)=u_0\in W^{1,\infty}(\R^N),
\end{split}
\end{equation}

\begin{equation}\label{condition h}
h\in \mathcal{C}^1(\R^{N+1}), |h(u,v)|\leq C(|u|^{\gamma}+|v|^{\overline{\gamma}}+1),
\end{equation}

\begin{equation}\label{behaviour1}
\left|\left|(T-t)^{\frac{1}{p-1}}u(.,t)-f\left(\frac{.-a}{\sqrt{(T-t)|\log(T-t)|}}\right)\right|\right|_{W^{1,\infty}(\R^N)}=O\left(\frac{1}{\sqrt{|\log(T-t)|}}\right),
\end{equation}

\begin{equation}\label{u_*}
u_*(x)=\underset{t\rightarrow T}{\lim}u(x,t)?
\end{equation}

\begin{equation}\label{region outside}
|x|\geq K_0\sqrt{(T_1-t)|\log(T_1-t)|)}.
\end{equation}

\begin{equation}\label{def t(x)}
|x|= \frac{K_0}{2}\sqrt{(T_1-t(x))|\log (T_1-t(x))|}.
\end{equation}

\begin{equation}\label{def: v}
\begin{split}
&v_1(\xi,\tau)=(T_1-t(x))^{\frac{1}{p-1}}u_1(x+\xi\sqrt{T_1-t(x)},t(x)+\tau (T_1-t(x))),\\
&v_2(\xi,\tau)=(T_1-t(x))^{\frac{1}{p-1}}\overset{\sim}{u}_2(x+\xi\sqrt{T_1-t(x)},t(x)+\tau (T_1-t(x))),
\end{split}
\end{equation}

\begin{equation}\label{profile v}
\begin{split}
\underset{\xi\in \R^N}{\sup}\left|(1-\tau)^{\frac{1}{p-1}}v(\xi,\tau)-f\left(z\right)\right|&\leq \frac{C}{\sqrt{|\log(T_1-t)|}}\\
&\leq \frac{C}{\sqrt{|\log(T_1-t(x))|}},
\end{split}
\end{equation}

\begin{equation}\label{z > K_0 beta}
|z|\geq \left(\frac{K_0}{2}-\frac{\beta}{|\log(T_1-t(x))|^{\frac{1}{4}}}\right).
\end{equation}

\begin{prp}[No blow up under a threshold]\label{prp bound v}
We assume that $v$ and $\nabla v$ satisfy the following equations:
\begin{equation}\label{ineq v}
\begin{split}
\partial_\tau v -\Delta v=f(v,\nabla v),
\end{split}
\end{equation}
\begin{equation}\label{ineq grad v}
\begin{split}
\partial_\tau \nabla v -\Delta \nabla v=\bar f(v,\nabla v),
\end{split}
\end{equation}
with $$|f(v,\nabla v)|\leq C(1+|v|^p+|\nabla v|^{\bar{\gamma}}),$$
$$|\bar f(v,\nabla v)-\nabla \bar g(v)|\leq C(1+|\nabla v|+|\nabla v||v|^{p-1}) \text{ and }|\bar g(v)|\leq C|\nabla v|^{\bar{\gamma}},$$
together with the following bounds, for all $ |\xi|<1$ and $\tau\in [0,1)$,
\begin{equation*}\label{bound v & nabla v}
|v(\xi,\tau)|+\sqrt{1-\tau}|\nabla v(\xi,\tau)|\leq \varepsilon_0(1-\tau)^{-\frac{1}{p-1}},
\end{equation*}
for some $\varepsilon_0>0$. Then, there exists $\varepsilon>0$, such that for all $|\xi|\leq \varepsilon,\tau \in [0,1)$,
$$|v(\xi,\tau)|+|\nabla v(\xi,\tau)|\leq C\varepsilon_0,$$
for some constant $C=C(p)>0$.
\end{prp}
One may see that the hypothesis of Proposition \ref{prp bound v} are verified in our case for small enough $|x|$. Therefore, we have the following direct consequence:
\begin{cor}\label{cor bound v, nabla v}
For any $\varepsilon_0 > 0$, there exist $K_0 > 0$ and $r_0 > 0$ such that, for all $x$ satisfying \eqref{region outside} and $|x| < r_0$, for all $|\xi|\leq\beta (|\log(T_1-t(x))|)^{\frac{1}{4}}-1,\tau \in [0,1)$, we have
\begin{equation}\label{bound v}
|v(\xi,\tau)|+|\nabla v(\xi,\tau)|\leq C\varepsilon_0.
\end{equation}
\end{cor}

\begin{equation}
|v(\xi,\tau)|+|\nabla v(\xi,\tau)|\leq C\varepsilon_0.
\end{equation}

\begin{equation}\label{chi a}
\chi=\chi_0\left(\frac{\xi}{a|\log(T_1-t(x))|^{\theta}}\right),
\end{equation}

\begin{equation}\label{equiv t(x)}
\log(T_1-t(x))\sim 2\log |x|\text{ and } T_1-t(x)\sim \frac{2|x|^2}{K_0^2|\log|x||}.
\end{equation}

\begin{theorem}[\textbf{Single-point blow-up and final profile}]\label{Thm 2 g}
Assume $N\geq 1$ and let be $(a,T)\in \R^N\times \R^+$, $h$ verifies \eqref{condition h} and $u\in S_{a,T,h}$.
\begin{enumerate}[label=(\roman*)]
\item It holds that $"a"$ is the only blow-up point.
\item There exists a final profile $u_*$ as in \eqref{u_*}, for all $x\in \R^N\backslash \{a\}$. Moreover, there exists $\epsilon>0$, such that for all $\epsilon\geq |x-a|> 0$\label{Thm 2 ii}
\begin{equation}\label{Thm 2}
\left|u_*(x)-\left(\frac{8p|\log|x-a||}{(p-1)^2|x-a|^{2}}\right)^{\frac{1}{p-1}}\right|\leq C\frac {|\log|x-a||^{\frac{1}{p-1} - \frac{1}{4}}} {|x-a|^{\frac{2}{p-1}}}.
\end{equation}
\end{enumerate}
\end{theorem}

\begin{proof}[Proof of Theorem \ref{Thm 2 g}]
Let us consider $u$ solution of \eqref{eq:1} such that $u\in S_{a,T,h}$. From equation \eqref{eq:1}, we may assume that $a=0$. Theorem \ref{eq:1} will follow if we prove the following two facts:
\begin{itemize}
\item The origin is the only blow-up point in $B(0,\epsilon)$ for some $\epsilon>0$, and $\ref{Thm 2 ii}$ holds in that same ball.
\item For any $|x|\geq \epsilon$, $x$ is not a blow-up point and $u(x,t)$ has a limit $u_*(x)$ as $t\rightarrow T$.
\end{itemize}
We proceed in two steps to prove these two facts.
\bigskip


\textbf{Step 1 (Single point blow-up in $B(0,\epsilon)$ and final profile):} One may see that the solution of the following equation: 
\begin{equation*}
\begin{split}
&v'_{K_0}=|v_{K_0}|^{p-1}v_{K_0},\\
&v_{K_0}(0)=f\left(\frac{K_0}{2}\right),
\end{split}
\end{equation*}
is given by the following: 
\begin{equation}\label{v_K_0}
v_{K_0}(\tau)=\left(f\left(\frac{K_0}{2}\right)^{1-p}-(p-1)\tau\right)^{-\frac{1}{p-1}}.
\end{equation}


We define $\eta(\xi,\tau)=v(\xi,\tau)-v_{K_0}(\tau)$, where $v$ is defined the same as $v_1$ in \eqref{def: v} with $T_1=T$ and $a=0$. We know from \eqref{bound v} and \eqref{v_K_0} that, for all $|\xi|\leq\beta (|\log(T_1-t(x))|)^{\frac{1}{4}}-1$ and $\tau \in [0,1)$, $|\eta(\xi,\tau)|<+\infty$ and $\eta$ verifies the following equation:
$$\partial_\tau \eta=\Delta \eta+\alpha_1 \eta+\alpha_2,$$
where 
\begin{equation*}
\begin{split}
&\alpha_1=\frac{|v|^{p-1}v-|v_{K_0}|^{p-1}v_{K_0}}{v-v_{K_0}},\\
&\alpha_2=(T_1-t(x))^{\frac{p}{p-1}}h\left((T_1-t(x))^{-\frac{1}{p-1}}v,(T-t(x))^{-\frac{p+1}{2(p-1)}}\nabla v\right).
\end{split}
\end{equation*}
Hence, we define again $\bar{\eta}=\eta \chi$, where $\chi$ is the same as in \eqref{chi a}. Thus, $\bar\eta$ verifies the following equation:
\begin{equation*}
\begin{split}
\partial_\tau \bar{\eta}=\chi \Delta \bar{\eta}-\frac{2\nabla(\eta\nabla \chi_0)}{a|\log (T-t(x))|^{\theta}}+\frac{\eta \Delta \chi_0}{a^2|\log (T-t(x))|^{2\theta}}+\alpha_1\bar{\eta}+\chi \alpha_2.\\
\end{split}
\end{equation*}
Therefore, it can be written as
\begin{equation*}
\begin{split}
\bar{\eta}(\tau)&=S(\tau) \bar{\eta}(0)-\frac{2}{a|\log (T_1-t(x))|^{\theta}}\int^\tau_0S(\tau-s)\nabla(\eta(s)\nabla \chi_0)ds\\
&+\frac{1}{a^2|\log (T_1-t(x))|^{2\theta}}\int^\tau_0S(\tau-s)\eta(s) \Delta \chi_0ds+\int^\tau_0 S(\tau-s)\alpha_1(s)\bar{\eta}(s)ds\\
&+\int^\tau_0 S(\tau-s)\chi\alpha_1(s)ds,
\end{split}
\end{equation*}
where $S$ is again the heat semigroup kernel. With the same computations as done in the previous section, we have


\begin{equation}\label{control bar eta}
\begin{split}
\|\bar{\eta}(\tau)\|_{L^\infty(\R^N)}&\leq\|\bar{\eta}(0)\|_{L^\infty(\R^N)}+\frac{C}{|\log (T_1-t(x))|^{\theta}}\int^\tau_0(\tau-s)^{-\frac{1}{2}}ds+\frac{C}{|\log (T_1-t(x))|^{2\theta}}\int^\tau_0ds\\
&+C\int^\tau_0 \|\bar{\eta}(s)\|_{L^\infty(\R^N)}ds+C(T_1-t(x))^{\nu}.
\end{split}
\end{equation}
Thus, for small $|x|$, we have
\begin{equation}\label{eta 0 left 2}
\begin{split}
\|\bar{\eta}\|_{L^\infty(\R^N)}(\tau)&\leq\|\bar{\eta}(0)\|_{L^\infty(\R^N)}+\frac{C}{|\log (T_1-t(x))|^{\theta}}+C\int^\tau_0 \|\bar{\eta}(s)\|_{L^\infty(\R^N)}ds.
\end{split}
\end{equation}
We know that $f$ is a Lipschitz function. Therefore,
\begin{equation}\label{f z-K_0}
\left|f(z)-f\left(\frac{K_0}{2}\right)\right|\leq C\left||z|-\frac{K_0}{2}\right|,
\end{equation}
where 
$$z=\frac{x+\xi\sqrt{T_1-t(x)}}{\sqrt{(T_1-t)|\log(T_1-t)|}}.$$
For $\tau =0$, we have $t=t(x)$, and for $x\geq 0$, we have, from \eqref{def t(x)},
\begin{equation*}
\begin{split}
|z|&\leq\frac{K_0}{2}\frac{\sqrt{(T_1-t(x))|\log(T_1-t(x))|}}{\sqrt{(T_1-t)|\log(T_1-t)|}}+\frac{|\xi|\sqrt{T_1-t(x)}}{\sqrt{(T_1-t)|\log(T_1-t)|}}\\
&=\frac{K_0}{2}+\frac{|\xi|}{\sqrt{|\log(T_1-t(x))|}}.\\
\end{split}
\end{equation*}
Since $|\xi|\leq \beta_1|\log(T_1-t(x))|^\theta$, then we get
%\begin{equation}
%|z-\frac{K_0}{2}|\leq \frac{1}{a|\log(T_1-t(x))|%^{\frac{1}{2}-\theta}}.
%\end{equation}
%With the same arguments for $x\leq 0$, we obtain 
\begin{equation*}
||z|-\frac{K_0}{2}|\leq \frac{1}{a|\log(T_1-t(x))|^{\frac{1}{2}-\theta}}.
\end{equation*}
Together with \eqref{profile v} and \eqref{f z-K_0}, we have 
\begin{equation*}
\underset{\xi\leq a |\log(T_1-t(x))|^{\theta}}{\sup}\left|v(\xi,0)-f(z)\right|\leq \frac{C}{|\log(T_1-t(x))|^{\frac{1}{2}-\theta}}.
\end{equation*}
Thus,
\begin{equation}\label{profile v(0)}
\underset{\xi\leq a |\log(T_1-t(x))|^{\theta}}{\sup}|\eta(\xi,0)|\leq \frac{C}{|\log(T_1-t(x))|^{\frac{1}{2}-\theta}}.
\end{equation}
Applying this to \eqref{eta 0 left 2}, we obtain
\begin{equation*}
\begin{split}
\|\bar{\eta}(\tau)\|_{L^\infty(\R^N)}&\leq \frac{C}{|\log(T_1-t(x))|^{\frac{1}{2}-\theta}}+\frac{C}{|\log (T_1-t(x))|^{\theta}}+C\int^\tau_0 \|\bar{\eta}(s)\|_{L^\infty(\R^N)}ds.
\end{split}
\end{equation*}
Using a Gronwall's argument again and by taking the optimal $\theta$ which is $\theta=\frac{1}{4}$, we get 
$$\|\bar{\eta}(\tau)\|_{L^\infty(\R^N)} \leq \frac{C}{|\log(T_1-t(x))|^{\frac{1}{4}}}.$$
Therefore, for $\xi\leq \beta_1|\log(T_1-t(x))|^{\frac{1}{4}}$ and $\tau \in [0,1)$,
$$|v(\xi,\tau)-v_{K_0}(\tau)|\leq \frac{C}{|\log(T_1-t(x))|^{\frac{1}{4}}}.$$
Then, for $\epsilon\in  (0; \frac{K_0}{2} \sqrt{T_1 \log T_1})$ and small $\delta_0$, for all $\epsilon\geq |x|\geq K_0 \sqrt{T_1-t \log (T_1-t)}$ and $t\in [T-\delta_0,T)$, we have
$$\left|u(x,t)-(T_1-t(x))^{-\frac{1}{p-1}}v_{K_0}\left(\frac{t-t(x)}{T_1-t(x)}\right)\right|\leq  C\frac{(T_1-t(x))^{-\frac{1}{p-1}}}{|\log(T_1-t(x))|^{\frac{1}{4}}}.$$
We define , for all $x\in \R\backslash \{a\}$ non-blow-up points, the following:
$$u_*(x)=\underset{t\rightarrow T}{\lim}u(x,t).$$
Then, for some  $\epsilon\in  (0; K_0 \sqrt{T \log T})$, we have for all $\epsilon\geq |x|\geq K_0 \sqrt{T-t \log (T-t)},$

$$|u_*(x)-(T-t(x))^{-\frac{1}{p-1}}v_{K_0}(1)|\leq  C\frac{(T_1-t(x))^{-\frac{1}{p-1}}}{|\log(T_1-t(x))|^{\frac{1}{4}}}.$$
Thus, from \eqref{equiv t(x)} with small enough $\delta_0$, we have for all $\epsilon\geq |x|\geq K_0 \sqrt{(T-t )\log (T-t)}$ and $t\in [T-\delta_0,T),$
\begin{equation}\label{final profile proof}
\left|u_*(x)-\left(\frac{2}{K_0^2}\right)^{-\frac{1}{p-1}}\frac{|x|^{-\frac{2}{p-1}}}{|\log|x||^{-\frac{1}{p-1}}}\left(f\left(\frac{K_0}{2}\right)^{1-p}-(p-1)\right)^{-\frac{1}{p-1}}\right|\leq  C(K_0) \frac {|\log|x-a||^{\frac{1}{p-1} - \frac{1}{4}}} {|x-a|^{\frac{2}{p-1}}},
\end{equation}
which gives us \eqref{Thm 2}. 

\textbf{Step 2 (No blow-up outside of $\mathbf{B(0,\epsilon)}$):} We may assume $T$ is as small as needed. Otherwise, we just have to consider in \eqref{eq:1}, $u_0(.)=u(.,T-\varepsilon)$, where $\varepsilon>0$ is small enough.

\textbf{Case of $\mathbf{\epsilon < |x|<\frac{K_0}{2}\sqrt{T|\log T|}}$:} We have from \eqref{z > K_0 beta}
$$|z|\geq \frac{K_0}{4}.$$
Together, with \eqref{profile v}, we obtain \begin{equation*}
(1-\tau)^{\frac{1}{p-1}}|v(\xi,\tau)|\leq \varepsilon_0, \forall |\xi|\leq\beta (|\log(T-t(x))|)^{\frac{1}{4}}, \forall \tau \in [0,1).
\end{equation*}
Easy similar computations yield to 
\begin{equation*}
|\nabla v(\xi,\tau)|\leq \varepsilon_0 (1-\tau)^{-\frac{1}{p-1}-\frac{1}{2}}, \forall \xi\in \R^N, \forall \tau \in [0,1).
\end{equation*}
With Proposition \ref{prp bound v}, we have for $|\xi|\leq\beta (|\log(T-t(x))|)^{\frac{1}{4}},\tau\in[0,1)$, 
$$|v(\xi,\tau)|\leq C.$$
Since $t(x)<T-\varepsilon$, for some $\varepsilon>0$, then 
\begin{equation}\label{control eta x>epsilon}
u_*(x)\leq C \text{ for all }\epsilon\leq |x|<\frac{K_0}{2}\sqrt{T|\log T|}.
\end{equation}

\textbf{Case of $\mathbf{|x| \geq \frac{K_0}{2}\sqrt{T|\log T|}}$:} We define for all $\tau(t)=\frac{t}{T}\in [0,1)$ and $|\xi|\leq \beta |\log T|^{\frac{1}{4}}$,
$$v(\xi,\tau)=T^{\frac{1}{p-1}}u(x+\xi\sqrt{T_1},\tau T).$$
Then, for $z=\frac{x+\xi\sqrt{T}}{\sqrt{(T-t)|\log(T-t)|}}$, we have
\begin{equation}\label{z K0}
\begin{split}
|z|&=\frac{\sqrt{T|\log T|}}{\sqrt{(T-t)|\log(T-t)|}}\left|\frac{x}{\sqrt{T|\log T|}}+\frac{\xi}{\sqrt{|\log T|}}\right|\\
&\geq \frac{\sqrt{T|\log T|}}{\sqrt{(T-t)|\log(T-t)|}}\left(\frac{K_0}{2}-\frac{\beta}{|\log T|^{\frac{1}{4}}}\right)\\
&\geq \frac{K_0}{4}.
\end{split}
\end{equation}
From \eqref{behaviour1}, we have
\begin{equation*}
\begin{split}
\underset{\xi\in \R^N}{\sup}\left|(1-\tau)^{\frac{1}{p-1}}v(\xi,\tau)-f\left(z\right)\right|&\leq \frac{C}{\sqrt{|\log(T-t)|}}\\
&\leq \frac{C}{\sqrt{|\log T|}}.
\end{split}
\end{equation*}
Together with \eqref{z K0}, we have for all $\tau(x,t)=\frac{t}{T}\in [0,1)$ and $|\xi|\leq \beta |\log T|^{\frac{1}{4}}$,
\begin{equation}\label{rough v 2}
(1-\tau)^{\frac{1}{p-1}}|v(\xi,\tau)|\leq \varepsilon_0.
\end{equation}
Again from \eqref{behaviour1}, we have
\begin{equation*}
\underset{x'\in \R^N}{\sup}\left|(1-\tau)^{\frac{1}{p-1}+\frac{1}{2}}\nabla v(\xi,\tau)-\frac{1}{\sqrt{|\log T|}}\nabla f\left(z\right)\right|\leq \frac{C}{\sqrt{|\log T|}},
\end{equation*}
which gives us 
\begin{equation*}
|\nabla v(\xi,\tau)|\leq \varepsilon_0 (1-\tau)^{-\frac{1}{p-1}-\frac{1}{2}}.
\end{equation*}
Together with \eqref{rough v 2} and Proposition \ref{prp bound v}, we obtain for $|\xi|\leq\beta (|\log(T-t(x))|)^{\frac{1}{4}},\tau\in[0,1)$, 
$$|v(\xi,\tau)|\leq C.$$
Again since $t(x)<T-\varepsilon$, then 

$$u_*(x)\leq C \text{ for all }|x| \geq \frac{K_0}{2}\sqrt{T|\log T|}.$$
Together with \eqref{final profile proof} and \eqref{control eta x>epsilon}, we conclude that $0$ is the only blow-up point and that $u_*(x)$ exists for all $x\in \R\backslash \{a\}$, which concludes the proof of Theorem \ref{Thm 2 g}.
\end{proof}

\end{document}
